\section{Coding sheet and coding manual}
\label{app:coding}

This supplement is the coding frame itself: all 38 dimensions in nine layers, each with its complete
permitted value set, reproduced from the frozen schema (v1.0.0) in the replication package. The main
text argues for the frame (\S5) and describes how it was applied (\S4);
neither argument is restated here. The schema was frozen on 2026-09-23, and no dimension has been added,
removed, renamed, or re-valued since the initial draft. Every table in this supplement is generated from
the schema by a script, so it cannot drift from what was coded.

The manuscript preamble loads \texttt{booktabs} and \texttt{graphicx} but not \texttt{longtable}, and a
section file cannot load a package, so the frame is given as nine per-layer tables
(Tables~\ref{tab:coding-A}--\ref{tab:coding-M}). A dimension is \emph{single} when one value is recorded
and \emph{multi} when every applicable value is recorded; 15 of the 38 are multi. Four dimensions
(\texttt{tool\_count}, \texttt{first\_release\_date}, \texttt{pinned\_version}, \texttt{stars}) carry a
scalar type instead of a value set. Each layer's table is preceded by the judgment call that layer
raised, illustrated by the two systems coded end to end in the manual,
SWE-agent v1.1.0~\citep{yang2024sweagent} and OpenHands v1.16.0~\citep{wang2024openhands}
(Section~\ref{subsec:coding-examples}).

\subsection{Coding sheet and judgment calls by layer}
\label{subsec:coding-sheet}

\paragraph{Layer A: context assembly.}\label{subsec:layer-a}
Layer A governs what the model sees at each step. A3 and A4 are multi-valued because compaction
mechanisms and observation modalities stack in one system, while a prompt-production strategy and a
context strategy do not. The judgment call on A1 is where a template stops. Variables rendered once per
run are \texttt{templated}; conditional sections, runtime-loaded files, or per-step recomposition are
\texttt{dynamically\_composed}; a paper that says only ``we prompt the model with'' is \texttt{static}
at low confidence. On A2, \texttt{mixed} is coded only when the default configuration both injects
context and lets the model navigate. The worked systems split on both dimensions. SWE-agent renders one
Jinja template per run and prunes tool output only: A1 \texttt{templated}, A3
\texttt{[tool\_output\_pruning]}. OpenHands composes a static tier from a section registry with a
dynamic tier from its agent context, and runs a summarizing condenser by default: A1
\texttt{dynamically\_composed}, A3 \texttt{[summarize, tool\_output\_pruning]}.

\begin{table}[htbp]
\caption{Layer \texttt{A}, Context assembly: 4 dimensions, with ids, names, keys, types and value
sets verbatim from the frozen schema (v1.0.0).}
\label{tab:coding-A}
\footnotesize
\begin{tabular}{@{}l p{0.20\linewidth} p{0.22\linewidth} l p{0.33\linewidth}@{}}
\toprule
ID & Dimension and key & What it records & Type & Permitted values \tabularnewline
\midrule
\texttt{A1} & \raggedright System-prompt style \newline \texttt{system\_prompt\_style} & \raggedright how the system prompt that opens every model call is produced & enum, single & \raggedright \texttt{static}, \texttt{templated}, \texttt{dynamically\_composed}, \texttt{model\_generated} \tabularnewline
\texttt{A2} & \raggedright Repo/environment context strategy \newline \texttt{env\_context\_strategy} & \raggedright how the model is given knowledge of the repository or environment & enum, single & \raggedright \texttt{none}, \texttt{file\_tree\_summary}, \texttt{retrieval\_bm25}, \texttt{embedding\_rag}, \texttt{agent\_driven\_navigation}, \texttt{mixed} \tabularnewline
\texttt{A3} & \raggedright Context compaction \newline \texttt{context\_compaction} & \raggedright mechanisms that shrink the history sent to the model & enum, multi & \raggedright \texttt{none}, \texttt{truncate\_oldest}, \texttt{summarize}, \texttt{tool\_output\_pruning}, \texttt{structured\_checkpoint}, \texttt{model\_native} \tabularnewline
\texttt{A4} & \raggedright Observation formatting \newline \texttt{observation\_format} & \raggedright the format in which tool results are presented to the model & enum, multi & \raggedright \texttt{raw\_text}, \texttt{structured\_json}, \texttt{screenshots}, \texttt{a11y\_tree}, \texttt{set\_of\_marks} \tabularnewline
\bottomrule
\end{tabular}
\end{table}

\paragraph{Layer B: tool interface.}\label{subsec:layer-b}
Layer B governs how actions are declared to the model and parsed back out of its output. Four of its
five dimensions are multi-valued, because harnesses ship several parsers, edit commands, and schema
sources at once; the exception, B2, is an integer and carries no value set. B2 is the layer's hardest
judgment call, because ``how many tools'' has no natural answer. The rule counts one entry per function
schema sent with the model call, and counts a built-in shell tool once regardless of what it can run. It
expands a toolset into its tools and counts built-ins attached to every agent, but excludes tools
injected by an MCP server or by the environment, and the maximal shipped configuration goes in the note.
SWE-agent is therefore \texttt{3} at the default against 17 tool definitions shipped across 13 bundles,
and OpenHands \texttt{5} against 19 once the 14-tool browser toolset is injected where Chromium exists.
Both counts are defensible and they are not the same measurement, which is why the counting rule rather
than the count is the taxonomy's content.

\begin{table}[htbp]
\caption{Layer \texttt{B}, Tool interface: 5 dimensions, with ids, names, keys, types and value
sets verbatim from the frozen schema (v1.0.0).}
\label{tab:coding-B}
\footnotesize
\begin{tabular}{@{}l p{0.20\linewidth} p{0.22\linewidth} l p{0.33\linewidth}@{}}
\toprule
ID & Dimension and key & What it records & Type & Permitted values \tabularnewline
\midrule
\texttt{B1} & \raggedright Tool-call format \newline \texttt{tool\_call\_format} & \raggedright the syntax by which the model requests an action & enum, multi & \raggedright \texttt{native\_function\_calling}, \texttt{json\_in\_text}, \texttt{xml\_tags}, \texttt{code\_as\_action}, \texttt{shell\_only}, \texttt{cli\_flags} \tabularnewline
\texttt{B2} & \raggedright Tool count at pinned version \newline \texttt{tool\_count} & \raggedright distinct tool definitions exposed to the model in the default configuration & integer, single & \raggedright no enumerated set (integer) \tabularnewline
\texttt{B3} & \raggedright Edit primitive \newline \texttt{edit\_primitive} & \raggedright how the model expresses a file modification & enum, multi & \raggedright \texttt{none}, \texttt{whole\_file\_rewrite}, \texttt{search\_replace}, \texttt{unified\_diff}, \texttt{line\_range\_edit}, \texttt{ast\_aware} \tabularnewline
\texttt{B4} & \raggedright Tool schema source \newline \texttt{tool\_schema\_source} & \raggedright where the tool schema shown to the model comes from & enum, multi & \raggedright \texttt{hand\_written}, \texttt{auto\_generated}, \texttt{mcp}, \texttt{openapi} \tabularnewline
\texttt{B5} & \raggedright Protocol standardization \newline \texttt{protocol\_standardization} & \raggedright standard agent or tool protocols the harness speaks & enum, multi & \raggedright \texttt{none}, \texttt{mcp}, \texttt{a2a}, \texttt{other} \tabularnewline
\bottomrule
\end{tabular}
\end{table}

\paragraph{Layer C: control loop.}\label{subsec:layer-c}
Layer C governs what happens between model calls: loop shape, plan representation, agent count,
hand-offs, and the points where a human can intervene. C1 and C5 are multi-valued because a harness can
run a ReAct loop inside an event stream and can gate at more than one point. The judgment call across the
whole layer is default versus capability. Decision rule~4 codes the shipped default and records
reachable alternatives in the note, so OpenHands, which ships a delegate tool but sets
\texttt{enable\_sub\_agents=False}, is C3 \texttt{single} and C4 \texttt{none}, with
\texttt{orchestrator\_workers} and \texttt{subagent\_spawn} named as opt-in. C2 draws a related line: a
structured plan the harness writes and re-reads across steps is \texttt{explicit\_plan\_object}, as in
OpenHands' task tracker, while a numbered recipe in the prompt is \texttt{implicit}, as in SWE-agent.

\begin{table}[htbp]
\caption{Layer \texttt{C}, Control loop: 5 dimensions, with ids, names, keys, types and value
sets verbatim from the frozen schema (v1.0.0).}
\label{tab:coding-C}
\footnotesize
\begin{tabular}{@{}l p{0.20\linewidth} p{0.22\linewidth} l p{0.33\linewidth}@{}}
\toprule
ID & Dimension and key & What it records & Type & Permitted values \tabularnewline
\midrule
\texttt{C1} & \raggedright Loop primitive(s) \newline \texttt{loop\_primitives} & \raggedright the control-flow patterns that drive model calls and actions & enum, multi & \raggedright \texttt{react}, \texttt{plan\_execute}, \texttt{generate\_test\_repair}, \texttt{multi\_attempt}, \texttt{tree\_search}, \texttt{event\_driven}, \texttt{fixed\_pipeline} \tabularnewline
\texttt{C2} & \raggedright Planning granularity \newline \texttt{planning\_granularity} & \raggedright whether and how a plan is represented & enum, single & \raggedright \texttt{none}, \texttt{implicit}, \texttt{explicit\_plan\_object}, \texttt{hierarchical} \tabularnewline
\texttt{C3} & \raggedright Multi-agent topology \newline \texttt{multi\_agent\_topology} & \raggedright the arrangement of agents in the default configuration & enum, single & \raggedright \texttt{single}, \texttt{orchestrator\_workers}, \texttt{peer}, \texttt{hierarchical}, \texttt{debate}, \texttt{pipeline} \tabularnewline
\texttt{C4} & \raggedright Delegation mechanism \newline \texttt{delegation\_mechanism} & \raggedright how work is handed to another agent, if at all & enum, single & \raggedright \texttt{none}, \texttt{subagent\_spawn}, \texttt{role\_handoff}, \texttt{message\_bus} \tabularnewline
\texttt{C5} & \raggedright Human-in-the-loop points \newline \texttt{human\_in\_loop} & \raggedright the points at which a human can gate or steer the run & enum, multi & \raggedright \texttt{none}, \texttt{on\_permission}, \texttt{on\_plan}, \texttt{on\_completion}, \texttt{configurable} \tabularnewline
\bottomrule
\end{tabular}
\end{table}

\paragraph{Layer D: memory and state.}\label{subsec:layer-d}
Layer D separates what outlives a model call from what outlives a run: working state inside a run (D1),
information carried between runs (D2), and whether an interrupted run can continue (D3). Only D2 is
multi-valued, because a harness can keep file notes and a skill library at once. The judgment call sits
on D3, where a trajectory file is not a checkpoint. Both \texttt{checkpoint} and \texttt{full\_resume}
require a code path that reads saved state back into a running agent. SWE-agent writes a per-step
trajectory but offers no path to resume from it, so it is \texttt{none} with the trajectory noted.
OpenHands reloads a conversation from a state file and an event directory, so it is
\texttt{full\_resume}. D2 excludes static few-shot demonstrations chosen in configuration, which are
prompt content rather than memory.

\begin{table}[htbp]
\caption{Layer \texttt{D}, Memory and state: 3 dimensions, with ids, names, keys, types and value
sets verbatim from the frozen schema (v1.0.0).}
\label{tab:coding-D}
\footnotesize
\begin{tabular}{@{}l p{0.20\linewidth} p{0.22\linewidth} l p{0.33\linewidth}@{}}
\toprule
ID & Dimension and key & What it records & Type & Permitted values \tabularnewline
\midrule
\texttt{D1} & \raggedright Short-term state \newline \texttt{short\_term\_state} & \raggedright what the harness keeps as working state within a run & enum, single & \raggedright \texttt{conversation\_only}, \texttt{scratchpad}, \texttt{structured\_task\_state} \tabularnewline
\texttt{D2} & \raggedright Long-term memory \newline \texttt{long\_term\_memory} & \raggedright information persisted across runs and re-injected & enum, multi & \raggedright \texttt{none}, \texttt{file\_notes}, \texttt{vector\_store}, \texttt{skill\_library}, \texttt{episodic\_db} \tabularnewline
\texttt{D3} & \raggedright State persistence across runs \newline \texttt{state\_persistence} & \raggedright whether an interrupted run can be continued & enum, single & \raggedright \texttt{none}, \texttt{checkpoint}, \texttt{full\_resume} \tabularnewline
\bottomrule
\end{tabular}
\end{table}

\paragraph{Layer E: verification and repair.}\label{subsec:layer-e}
Layer E records checks the harness enforces, not checks the model could choose to perform. E1 is
multi-valued because linting, test execution, and a judge model coexist, while a retry policy and a
rollback mechanism are one choice each. The enforcement boundary is the layer's main judgment call: if
the model may run tests through a shell but the harness neither runs them nor reads their result,
\texttt{test\_execution} is not coded. E3 shows the closest-value rule at work. Both worked systems
expose a model-invocable per-file \texttt{undo\_edit} backed by saved copies, which is not a filesystem
or container snapshot; both are coded \texttt{snapshot} with the mismatch written into the note and
listed among the ambiguities (Table~\ref{tab:ambiguities-full}).

\begin{table}[htbp]
\caption{Layer \texttt{E}, Verification and repair: 3 dimensions, with ids, names, keys, types and value
sets verbatim from the frozen schema (v1.0.0).}
\label{tab:coding-E}
\footnotesize
\begin{tabular}{@{}l p{0.20\linewidth} p{0.22\linewidth} l p{0.33\linewidth}@{}}
\toprule
ID & Dimension and key & What it records & Type & Permitted values \tabularnewline
\midrule
\texttt{E1} & \raggedright Self-verification \newline \texttt{self\_verification} & \raggedright checks the harness applies, or prompts for, on the agent's work & enum, multi & \raggedright \texttt{none}, \texttt{self\_critique}, \texttt{test\_execution}, \texttt{linters\_typecheck}, \texttt{llm\_judge}, \texttt{formal} \tabularnewline
\texttt{E2} & \raggedright Retry policy \newline \texttt{retry\_policy} & \raggedright task-level retry behavior after a failed or rejected attempt & enum, single & \raggedright \texttt{none}, \texttt{fixed\_n}, \texttt{until\_pass}, \texttt{adaptive} \tabularnewline
\texttt{E3} & \raggedright Rollback / undo \newline \texttt{rollback} & \raggedright the mechanism that reverts changes made during the run & enum, single & \raggedright \texttt{none}, \texttt{git\_based}, \texttt{snapshot} \tabularnewline
\bottomrule
\end{tabular}
\end{table}

\paragraph{Layer F: budget and termination.}\label{subsec:layer-f}
Layer F records what bounds a run. Termination conditions, cost controls, and timeouts each admit a
set, because bounds accumulate rather than exclude one another. A condition counts only if it stops the
loop rather than merely logging, and an off-by-default cap counts when a configuration field exists,
with the default stated in the note. The layer's ambiguity is what \texttt{max\_steps} and
\texttt{max\_tokens} mean in practice: SWE-agent caps API calls rather than agent steps, and exits on
context-window overflow rather than on a token budget. F3 is narrower than it looks, because iteration
caps are not timeouts. SWE-agent is \texttt{[per\_tool, per\_run]}, where the per-run limit counts
cumulative command execution time and not wall-clock; OpenHands is \texttt{[per\_tool]}, with the
absence of a run-level timeout evidenced by the search that found none.

\begin{table}[htbp]
\caption{Layer \texttt{F}, Budget and termination: 3 dimensions, with ids, names, keys, types and value
sets verbatim from the frozen schema (v1.0.0).}
\label{tab:coding-F}
\footnotesize
\begin{tabular}{@{}l p{0.20\linewidth} p{0.22\linewidth} l p{0.33\linewidth}@{}}
\toprule
ID & Dimension and key & What it records & Type & Permitted values \tabularnewline
\midrule
\texttt{F1} & \raggedright Termination condition \newline \texttt{termination\_condition} & \raggedright the conditions that end a run & enum, multi & \raggedright \texttt{model\_declared}, \texttt{max\_steps}, \texttt{max\_tokens}, \texttt{cost\_cap}, \texttt{test\_pass}, \texttt{stall\_detection} \tabularnewline
\texttt{F2} & \raggedright Cost controls \newline \texttt{cost\_controls} & \raggedright mechanisms that limit or reduce spend & enum, multi & \raggedright \texttt{none}, \texttt{token\_budget}, \texttt{model\_routing}, \texttt{caching} \tabularnewline
\texttt{F3} & \raggedright Timeouts \newline \texttt{timeouts} & \raggedright time limits the harness enforces & enum, multi & \raggedright \texttt{none}, \texttt{per\_tool}, \texttt{per\_run} \tabularnewline
\bottomrule
\end{tabular}
\end{table}

\paragraph{Layer G: sandbox and environment.}\label{subsec:layer-g}
Layer G records where actions execute and what they may touch. It is the only layer whose every
dimension is single-valued: isolation boundary, filesystem scope, egress policy, and authorization model
are each one configured choice at the default. G2 carries a reference-frame problem that the taxonomy
resolves by convention rather than by adding values. The value \texttt{full} is coded relative to the
G1 boundary, so SWE-agent's root access inside a Docker container and OpenHands' documented full access
to the host filesystem are both \texttt{full}, with the note stating which frame applies. G2 must
therefore be read together with G1 and never alone. G3 is where the absence rule
(\S5 and decision rule~5 below) does most of its work. Neither worked system implements an egress policy, so
each is \texttt{open} at medium confidence, evidenced by the workspace configuration that would declare
a policy and by the search that found none. In G3, unlike every other dimension, \texttt{none} is not
the absence value: it means no network access, and the absence of any egress policy is \texttt{open}.

\begin{table}[htbp]
\caption{Layer \texttt{G}, Sandbox and environment: 4 dimensions, with ids, names, keys, types and value
sets verbatim from the frozen schema (v1.0.0).}
\label{tab:coding-G}
\footnotesize
\begin{tabular}{@{}l p{0.20\linewidth} p{0.22\linewidth} l p{0.33\linewidth}@{}}
\toprule
ID & Dimension and key & What it records & Type & Permitted values \tabularnewline
\midrule
\texttt{G1} & \raggedright Execution isolation \newline \texttt{execution\_isolation} & \raggedright the boundary within which agent actions execute, at the default & enum, single & \raggedright \texttt{none}, \texttt{subprocess}, \texttt{container}, \texttt{vm}, \texttt{remote} \tabularnewline
\texttt{G2} & \raggedright Filesystem access \newline \texttt{filesystem\_access} & \raggedright what the agent may read or write inside the G1 boundary & enum, single & \raggedright \texttt{full}, \texttt{scoped}, \texttt{read\_only}, \texttt{virtual} \tabularnewline
\texttt{G3} & \raggedright Network policy \newline \texttt{network\_policy} & \raggedright network egress available to agent actions & enum, single & \raggedright \texttt{open}, \texttt{allowlist}, \texttt{none} \tabularnewline
\texttt{G4} & \raggedright Permission model \newline \texttt{permission\_model} & \raggedright how individual actions are authorized & enum, single & \raggedright \texttt{none}, \texttt{static\_allowlist}, \texttt{per\_call\_prompt}, \texttt{policy\_engine} \tabularnewline
\bottomrule
\end{tabular}
\end{table}

\paragraph{Layer H: observability and governance.}\label{subsec:layer-h}
Layer H records what a run leaves behind and what constrains it, and it keeps observability and
governance together because in this corpus they are reported together or not at all. Only H4 is
multi-valued. H1 is coded at the highest level the pinned artifact can reach, even when that level needs
environment variables, and the note states what runs by default. OpenHands is \texttt{opentelemetry} on
the strength of an OTLP exporter, although only its structured event log runs unconditionally. H2
refuses \texttt{deterministic} without an explicit statement, so both worked systems are
\texttt{partial}. H3 is scoped to the pinned repositories. OpenHands is \texttt{none} although its 2024
paper described an integrated evaluation framework, because at the coded version that code lives outside
those repositories; the paper's claim goes in the note.

\begin{table}[htbp]
\caption{Layer \texttt{H}, Observability and governance: 4 dimensions, with ids, names, keys, types and value
sets verbatim from the frozen schema (v1.0.0).}
\label{tab:coding-H}
\footnotesize
\begin{tabular}{@{}l p{0.20\linewidth} p{0.22\linewidth} l p{0.33\linewidth}@{}}
\toprule
ID & Dimension and key & What it records & Type & Permitted values \tabularnewline
\midrule
\texttt{H1} & \raggedright Tracing \newline \texttt{tracing} & \raggedright the richest form of run record the harness can emit & enum, single & \raggedright \texttt{none}, \texttt{logs}, \texttt{structured\_traces}, \texttt{opentelemetry} \tabularnewline
\texttt{H2} & \raggedright Replayability \newline \texttt{replayability} & \raggedright whether a recorded trajectory can be re-run & enum, single & \raggedright \texttt{none}, \texttt{partial}, \texttt{deterministic} \tabularnewline
\texttt{H3} & \raggedright Evaluation hooks \newline \texttt{eval\_hooks} & \raggedright benchmark evaluation wired into the harness's own repositories & enum, single & \raggedright \texttt{none}, \texttt{built\_in} \tabularnewline
\texttt{H4} & \raggedright Guardrails \newline \texttt{guardrails} & \raggedright safety controls on inputs, outputs, and actions & enum, multi & \raggedright \texttt{none}, \texttt{input\_filters}, \texttt{output\_filters}, \texttt{action\_policies} \tabularnewline
\bottomrule
\end{tabular}
\end{table}

\paragraph{Layer M: meta.}\label{subsec:layer-m}
Layer M is metadata on the system rather than a clause of the definition, and it exists so that the
corpus can be stratified, weighted, and ordered in time. Three of its seven dimensions carry no value set
because they are a date, a string, and an integer. M6 is the dimension the rest of the coding depends on,
since every evidence string is relative to it. Multi-repository systems are why it is a string rather
than an enumeration. OpenHands is pinned at both its flagship release and the agent SDK release that the
flagship pins, because the harness now lives in the second repository while the name belongs to the
first. M4 is \texttt{repo} whenever the pinned repository diverges from the paper, as in both worked
systems, with the gap noted.

\begin{table}[htbp]
\caption{Layer \texttt{M}, Meta: 7 dimensions, with ids, names, keys, types and value
sets verbatim from the frozen schema (v1.0.0).}
\label{tab:coding-M}
\footnotesize
\begin{tabular}{@{}l p{0.20\linewidth} p{0.22\linewidth} l p{0.33\linewidth}@{}}
\toprule
ID & Dimension and key & What it records & Type & Permitted values \tabularnewline
\midrule
\texttt{M1} & \raggedright Target domain \newline \texttt{target\_domain} & \raggedright the task domains the harness is built for & enum, multi & \raggedright \texttt{swe}, \texttt{gui\_computer\_use}, \texttt{web}, \texttt{general\_tool\_use}, \texttt{research}, \texttt{other} \tabularnewline
\texttt{M2} & \raggedright Open source \newline \texttt{open\_source} & \raggedright whether the source is open & enum, single & \raggedright \texttt{yes}, \texttt{no}, \texttt{partial} \tabularnewline
\texttt{M3} & \raggedright Model agnostic \newline \texttt{model\_agnostic} & \raggedright whether the harness runs across model providers & enum, single & \raggedright \texttt{yes}, \texttt{no} \tabularnewline
\texttt{M4} & \raggedright Primary artifact \newline \texttt{primary\_artifact} & \raggedright the artifact the pinned version is coded from & enum, single & \raggedright \texttt{paper}, \texttt{repo}, \texttt{tech\_report} \tabularnewline
\texttt{M5} & \raggedright First release date \newline \texttt{first\_release\_date} & \raggedright first tagged release of the coded line & date, single & \raggedright no enumerated set (date) \tabularnewline
\texttt{M6} & \raggedright Pinned commit or version \newline \texttt{pinned\_version} & \raggedright tag and full commit hash for every repository coded & string, single & \raggedright no enumerated set (string) \tabularnewline
\texttt{M7} & \raggedright GitHub stars (date-stamped in evidence) \newline \texttt{stars} & \raggedright GitHub stars of the flagship repository, date-stamped in evidence & integer, single & \raggedright no enumerated set (integer) \tabularnewline
\bottomrule
\end{tabular}
\end{table}

\subsection{Decision rules}
\label{subsec:coding-rules}

The coding manual in the replication package (v1.0, frozen 2026-09-23) states nine general rules. They
are summarized here in the order the manual gives them; rule~5 is reproduced in full because the
under-reporting result depends on it.

\begin{enumerate}
\item \textbf{Pin first.} Code the system at its pinned version (\texttt{M6}). Where the paper and
  the repository disagree, code the repository at the pinned commit and record the paper's claim in
  \texttt{note}. A harness spanning several repositories has each pinned, with evidence paths
  prefixed by the repository name.
\item \textbf{Evidence is verbatim.} Either a repository locator or a verbatim paper quote with its
  section number; no paraphrase. A coder who cannot point at evidence records
  \texttt{not\_reported: true} and \texttt{value: null}, except where the value being recorded
  \emph{is} an absence, which rule~5 governs.
\item \textbf{Confidence.} \texttt{high} for an explicit statement or code on the cited line,
  \texttt{medium} for an inference from adjacent code, a default, or a figure, \texttt{low} for an
  inference from prose describing behavior.
\item \textbf{Default first, capability second.} Code the shipped default configuration of the
  pinned artifact. For multi-valued dimensions, add every value reachable through shipped
  configuration without code changes and name the default in \texttt{note}; for single-valued
  dimensions, code the default and list the alternatives in \texttt{note}. Never record
  \texttt{none} alongside another value.
\item \textbf{Absence is not silence.} ``I looked and found nothing'' splits into two codings, and
  the split decides the under-reporting rate, so it is decided in this order:
  \begin{enumerate}
  \item Does the evidence contain \emph{the place where this feature would be declared if it
    existed}: the config schema or settings model, the CLI flags, the tool or plugin registry,
    the run loop itself, or a feature list in the documentation?
  \item \textbf{Yes, and the feature is not there} $\rightarrow$ code the absence value
    (\texttt{none}, \texttt{open}, and so on). The evidence is that place: the config model,
    registry, or loop that was opened, plus the search that was run. Confidence is at most
    \texttt{medium} unless the documentation states the absence outright.
  \item \textbf{No, the evidence has no such place} (no config surface in the bundle, a truncated
    repository, or a paper only) $\rightarrow$ \texttt{not\_reported: true}, \texttt{value: null},
    and \texttt{note} names the place that was missing.
  \end{enumerate}
  The absence of a \emph{mention} is never evidence of absence: prose that does not discuss rollback
  is not evidence that rollback is absent. Equally, \texttt{not\_reported} is never correct once the
  place where the feature would be declared has been opened. This is the form rule~5 took under
  protocol amendment~9; the earlier form and the measured effect of the rewrite are reported in
  \S4.
\item \textbf{Closest value plus note.} Where a value list does not fit, record the closest value,
  write the mismatch in \texttt{note}, and add the case to the manual's ambiguity list, which holds
  26 entries, all listed with their resolutions in Table~\ref{tab:ambiguities-full}. Value lists change only through the schema plus a
  changelog line.
\item \textbf{Line numbers are real.} Cite the line read at the pinned commit; ranges are allowed.
  Amendment~8 relaxed the locator to \texttt{path@commit} where the evidence bundle presents
  repository excerpts without their original line numbers, and flags every such cell
  \texttt{locator\_no\_line}.
\item \textbf{Paper quotes} carry the arXiv identifier and the section.
\item \textbf{Coder id} records who or what produced the cell.
\end{enumerate}

The ordered test in rule~5 is what makes \texttt{not\_reported} a measurement claim about the
sources rather than a claim about the system, and the two states are reported separately throughout:
a cell is a value, or \texttt{not\_reported} (the sources were read and are silent), or
\texttt{unresolved} (the coder could not settle it). Counts for all three are in
Section~\ref{app:systems}. The sources are silent most often on \texttt{replayability} unweighted, at
85.1\% of its 1,256 coded cells, and on \texttt{network\_policy} weighted to the census, at 93.3\%
(84.8\% of its 1,253 resolved cells unweighted); the per-dimension rates are released in the replication
package.

Two boundary rules belong to the frame rather than to any one dimension. A framework is coded only
through its shipped default agent, at the default configuration, with the file defining that default as
evidence, so framework rows rest on a stated convention rather than a hidden judgment. A benchmark's
reference agent is a coded system when the benchmark ships a looping baseline, and the benchmark itself
never is.

Multi-valued cells need a reliability reading beside the registered one. Exact-set matching scores two
readings that agree on six of seven primitives as a complete disagreement, so its penalty grows with the
size of the true set rather than with the size of the error. On the 247-system double-coded sample,
\texttt{loop\_primitives} scores 0.590 on the registered exact-set Cohen's kappa, 0.628 on Gwet's AC1,
and 0.731 on per-value kappa, with per-value agreement of 0.928 across its seven values. All four are
measurements of reproducibility between two model readings, and none measures correctness. The
dimension is kept and flagged as the weakest rather than redefined.

\subsection{Worked examples}
\label{subsec:coding-examples}

The manual carries two systems coded end to end on every dimension, released in the replication package
as its two worked-example files. Both are LLM pre-fill and are not human-verified, which amendment~4
made the rule for the whole dataset rather than an exception for these two; both carry 38 cells with no
\texttt{not\_reported}. SWE-agent~\citep{yang2024sweagent} is pinned at \texttt{v1.1.0} (commit
\texttt{0f3acaf}, 2025-05-22), so its locators are repository-root relative.
OpenHands~\citep{wang2024openhands} is pinned at \texttt{v1.16.0} (commit \texttt{64c126965},
2026-08-27), which pins \texttt{software-agent-sdk v1.44.0} (commit \texttt{322dec7d}). At that tag the
flagship repository is a frontend and the agent loop lives in the SDK, so agent-behavior cells cite the
SDK and every locator carries its repository name as a prefix. The pair is the reason the manual's
ambiguity list exists: each of the 26 entries names a case these two systems raised.

Table~\ref{tab:coding-examples} gives eight of the 38 cells from each, chosen to span the layers; it
shows how the two systems differ on the same frame. The full cells, with evidence string, confidence,
and note, are in the two released files.

\begin{table}[htbp]
\caption{Eight of the 38 cells of each worked example, verbatim from the released worked-example
files. Values only; the evidence, confidence, and note of every cell are in those files.}
\label{tab:coding-examples}
\footnotesize
\begin{tabular}{@{}l l p{0.28\linewidth} p{0.30\linewidth}@{}}
\toprule
ID & key & SWE-agent 1.x & OpenHands \tabularnewline
\midrule
\texttt{A1} & \texttt{system\_prompt\_style} & \raggedright \texttt{templated} & \raggedright \texttt{dynamically\_composed} \tabularnewline
\texttt{C1} & \texttt{loop\_primitives} & \raggedright \texttt{react}, \texttt{multi\_attempt} & \raggedright \texttt{react}, \texttt{event\_driven}, \texttt{generate\_test\_repair} \tabularnewline
\texttt{B1} & \texttt{tool\_call\_format} & \raggedright \texttt{native\_function\_calling}, \texttt{json\_in\_text}, \texttt{xml\_tags}, \texttt{shell\_only} & \raggedright \texttt{native\_function\_calling}, \texttt{xml\_tags} \tabularnewline
\texttt{B2} & \texttt{tool\_count} & \raggedright \texttt{3} & \raggedright \texttt{5} \tabularnewline
\texttt{E1} & \texttt{self\_verification} & \raggedright \texttt{self\_critique}, \texttt{linters\_typecheck}, \texttt{llm\_judge} & \raggedright \texttt{llm\_judge} \tabularnewline
\texttt{G1} & \texttt{execution\_isolation} & \raggedright \texttt{container} & \raggedright \texttt{subprocess} \tabularnewline
\texttt{G3} & \texttt{network\_policy} & \raggedright \texttt{open} & \raggedright \texttt{open} \tabularnewline
\texttt{H2} & \texttt{replayability} & \raggedright \texttt{partial} & \raggedright \texttt{partial} \tabularnewline
\bottomrule
\end{tabular}
\end{table}

Table~\ref{tab:ambiguities-full} lists all 26 entries of the manual's ambiguity list with the rule
that resolves each.

\begin{table}[htbp]
\caption{The 26 ambiguities recorded while coding the two worked systems, each resolved by a
coding-manual rule and none by changing a value set. Numbering follows the manual; \S5.3 of the main
paper discusses six of them.}
\label{tab:ambiguities-full}
\footnotesize
\begin{tabular}{@{}r p{0.17\linewidth} p{0.31\linewidth} p{0.40\linewidth}@{}}
\toprule
\# & Dimension & Ambiguity & Resolution \tabularnewline
\midrule
1 & B2, C3--C5, D3, E1, E2, G1, G4 & \raggedright Capability shipped but off by default & \raggedright Code the default; name reachable alternatives in the note \tabularnewline
2 & B2 & \raggedright Toolsets, built-in, injected, and MCP tools & \raggedright One per schema sent; toolsets expanded; injected tools excluded \tabularnewline
3 & B3 & \raggedright New-file creation; custom patch formats & \raggedright New file \texttt{whole\_file\_rewrite}; other patches \texttt{unified\_diff}, syntax named \tabularnewline
4 & B4 & \raggedright Hand-written declarative specs converted to schema & \raggedright \texttt{hand\_written}: human-authored in any format; \texttt{auto\_generated}: derived from code \tabularnewline
5 & B1 & \raggedright Command in a fenced block & \raggedright Fenced command \texttt{shell\_only}; tag-delimited call \texttt{xml\_tags} \tabularnewline
6 & C1 & \raggedright Refinement driven by a judge model & \raggedright \texttt{generate\_test\_repair}, verifier type noted \tabularnewline
7 & C5 & \raggedright A human replaces the model & \raggedright \texttt{none}; \texttt{configurable} means runtime-selectable gating policy \tabularnewline
8 & D3 & \raggedright Trajectory written but not resumable & \raggedright \texttt{none}; trace writing alone is not a checkpoint \tabularnewline
9 & E2 & \raggedright Step requery versus task-level retry & \raggedright Task-level only; step and transport retries noted \tabularnewline
10 & E3 & \raggedright Per-file tool-level undo & \raggedright \texttt{snapshot}; a repository reset is \texttt{git\_based} \tabularnewline
11 & F1 & \raggedright Caps on API calls; exit on context overflow & \raggedright \texttt{max\_steps}: iterations or calls; \texttt{max\_tokens}: budget or overflow \tabularnewline
12 & F2 & \raggedright Budget capped in currency & \raggedright \texttt{token\_budget}, currency noted \tabularnewline
13 & G1 & \raggedright Several isolation backends shipped & \raggedright Default coded, others noted \tabularnewline
14 & G2 & \raggedright Full access in a container versus on the host & \raggedright \texttt{full} relative to G1; note states the frame \tabularnewline
15 & G3 & \raggedright No network configuration anywhere & \raggedright \texttt{open} at medium confidence, citing the configuration surface \tabularnewline
16 & G4 & \raggedright Static blocklist & \raggedright \texttt{static\_allowlist}, noted as a blocklist \tabularnewline
17 & H1 & \raggedright Optional trace exporter versus always-on log & \raggedright Highest available level; default noted \tabularnewline
18 & H3 & \raggedright Benchmarks in a sibling repository & \raggedright \texttt{built\_in} only inside the pinned repositories \tabularnewline
19 & M4 & \raggedright Paper describes an older version & \raggedright \texttt{repo}, with the version gap noted \tabularnewline
20 & M5 & \raggedright Which release counts as first & \raggedright First tag of coded line, then first commit, then paper \tabularnewline
21 & M6 & \raggedright Harness spans several repositories & \raggedright Tag, commit, and date per repository \tabularnewline
22 & M7 & \raggedright Stars across several repositories & \raggedright Flagship repository counted; others noted \tabularnewline
23 & A2 & \raggedright Repository instruction files inject context & \raggedright Noted; no new value added \tabularnewline
24 & A4 & \raggedright DOM element list rendered as text & \raggedright \texttt{a11y\_tree} \tabularnewline
25 & B5 & \raggedright Agent Client Protocol & \raggedright \texttt{other}, protocol named; no \texttt{acp} value added \tabularnewline
26 & General & \raggedright Evidence for an absence value & \raggedright The absence rule (\S5.2 of the main paper) \tabularnewline
\bottomrule
\end{tabular}
\end{table}
